% Assignment 5 (Q1) - Report
% Tarun Konde (2025UCP1811)
% The Beamer part (Q2) is in ../../Beamer/presentation.tex

\documentclass[12pt,a4paper]{report}

\usepackage[T1]{fontenc}
\usepackage{lmodern}
\usepackage[margin=1in]{geometry}
\usepackage[dvipsnames]{xcolor}
\usepackage{sectsty}
\usepackage{fancyhdr}
\usepackage{amsmath}

% (c) heading styles with sectsty
\chapterfont{\color{red}\bfseries\huge}
\sectionfont{\color{magenta}\itshape\large}
\subsectionfont{\color{orange}}

% (d) header from the second page onwards
\pagestyle{fancy}
\fancyhf{}
\fancyhead[L]{Tarun Konde}
\fancyhead[C]{Technical Writing}
\fancyhead[R]{\today}
\fancyfoot[C]{\thepage}
% the report class switches chapter opening pages to "plain",
% so plain is redefined to carry the same header
\fancypagestyle{plain}{%
    \fancyhf{}
    \fancyhead[L]{Tarun Konde}
    \fancyhead[C]{Technical Writing}
    \fancyhead[R]{\today}
    \fancyfoot[C]{\thepage}
}

\begin{document}

% (a) title page - no header or footer here
\begin{titlepage}
    \thispagestyle{empty}
    \centering
    \vspace*{3cm}
    {\Huge\bfseries Introduction to \LaTeX\par}
    \vspace{2cm}
    {\Large Tarun Konde\par}
    \vspace{0.4cm}
    {\large 2025UCP1811\par}
    \vspace{2cm}
    {\large Course: Technical Writing (22CSP210)\par}
    \vspace{0.3cm}
    {\large Instructor: Prof.\ Neeta Nain\par}
    \vfill
    {\large \today\par}
\end{titlepage}

% (b) abstract
% The abstract environment of the report class prints on its own
% title-style page with no header, so a starred chapter is used instead.
\chapter*{Abstract}
This report is a short introduction to \LaTeX{} written as part of the
Technical Writing lab. It explains what \LaTeX{} is, how a document is
organised, and how mathematics, tables and figures are handled. It also lists
the problems I ran into while learning it.

\chapter{Getting Started}

\section{Introduction}

\LaTeX{} is a document preparation system built on top of \TeX, which was
written by Donald Knuth. The author writes plain text with markup commands, and
the compiler produces the final PDF. The separation of content and formatting
is the main difference from Word or Google Docs.

\section{Installing the tools}

I used TeX Live with TeXstudio as the editor. Overleaf works the same way in
the browser and needs no installation.

\subsection{Compiling}

A document is compiled with \texttt{pdflatex}. When the document has
references or a bibliography it has to be compiled more than once, because the
first run only collects the labels.

\chapter{Structure of a Document}

\section{Preamble and body}

Everything before \verb|\begin{document}| is the preamble. Packages and
settings go there. The text goes in the body.

\section{Sectioning}

The \texttt{report} class has chapters, sections, subsections and
paragraphs. Numbering and the table of contents are generated from these
commands.

\subsection{Headers and footers}

The \texttt{fancyhdr} package controls headers and footers. In this report the
title page has none, and every page after it shows my name, the course and the
date.

\chapter{Mathematics}

\section{Inline and display math}

Inline math is written between dollar signs, for example $a^2 + b^2 = c^2$.
Display math gets its own line:
\begin{equation}
    \int_0^1 x^2\,dx = \frac{1}{3}
\end{equation}

\subsection{What I found difficult}

Matching braces in long equations and remembering which package defines a
command took the most time at first.

\chapter{Conclusion}

\section{Summary}

\LaTeX{} takes longer to learn than a word processor, but once the basics are
clear it saves time on long documents with equations and references.

\end{document}
